%% 第三章

\chapter{B-rep 序列化与数据预处理}
\chaptermark{B-rep 序列化与数据预处理}

本章介绍把 B-rep 模型转换为可训练的 token 序列的整套预处理流程，主要内容包括：数据流程总览、B-rep 几何 / 拓扑分解、面 / 边 NURBS 采样、坐标量化、统一词表、序列骨架（face\_block / edge\_block）的组织方式，以及面 / 边的拓扑感知排序策略。

\section{数据预处理流程总览}

整套预处理从 STEP 文件出发，经五个阶段抵达可训练的序列文件：STEP 解析 $\to$ 拓扑分解 $\to$ 面 / 边几何采样 $\to$ 坐标量化 $\to$ 拓扑感知排序与序列拼装。其中 STEP 解析与拓扑分解通过 OpenCASCADE 完成。

本文实现中保留了原始 NURBS 采样版、量化版以及服务器侧镜像脚本，三者在核心采样、量化与排序规则上保持一致。最终训练数据包含 64\,325 个训练样本与 6\,095 个验证样本，统一词表规模 4\,053；具体脚本、数据文件和日志路径统一列入附录，正文不再逐一展开。

\begin{figure}[!htbp]
    \centering
    \includegraphics[width=0.82\textwidth]{插图/fig3-1_preprocess_pipeline.pdf}
    \caption{数据预处理流程示意图。流程以 STEP 文件中的 B-rep 输入为起点，依次经过 OpenCASCADE 解析、拓扑分解、面 / 边 NURBS 采样、坐标归一化、4000 级标量量化和拓扑感知排序，最终拼装为包含 START、SEP、END 等结构 token 的训练序列数据集。}
    \label{fig:preprocess-pipeline}
\end{figure}

\section{B-rep 的几何与拓扑分解}

\subsection{面 / 边 / 顶点的提取}

读入一个 STEP 文件后，使用 OpenCASCADE 的 \texttt{TopExp\_Explorer} 按 Face、Edge、Vertex 顺序遍历拓扑层级，得到面集合 $\{f_k\}_{k=1}^{F}$ 与边集合 $\{e_j\}_{j=1}^{E}$，并维护面—边邻接表：对每个面记录其所有 Wire 上的边索引，对每个边记录其相邻的两个面索引。这一阶段不修改任何几何，只构建索引表与邻接图。

\subsection{NURBS 曲面的 UV 采样}

对每个面 $f_k$，在其参数域 $[u_{\min}, u_{\max}] \times [v_{\min}, v_{\max}]$ 上做均匀网格采样，得到 \textbf{16 个三维点}，作为该面在序列中的几何样本：

\begin{equation}
    \mathcal{P}_k = \{ S_k(u_a, v_b) \mid a, b = 0, 1, 2, 3 \}.
    \label{eq:face-uv-sampling}
\end{equation}

为减小参数化任意性带来的分布偏差，预处理阶段对每个面 UV 方向做 canonical sorting：按面包围盒主轴对 $u$、$v$ 方向重新排序。该 16 个三维点共贡献 48 个标量坐标 token，与 1 个面索引 token 一起构成 \textbf{face\_block = 49} 的面段（详见 §3.4）。

\subsection{NURBS 曲线的弧长采样}

对每条边 $e_j$，按弧长在 $[0, L_j]$ 上均匀取参数点并求值，得到 \textbf{4 个三维点} 作为该边在序列中的几何样本。弧长采样比参数域均匀采样更稳健——对参数化变形不敏感，对短边和长边都能给出可比较的采样密度。Edge 同样做 canonical sorting：按起止顶点的坐标顺序统一方向。该 4 个三维点共贡献 12 个标量坐标 token，与 2 个面索引 token（两端邻面）一起构成 \textbf{edge\_block = 14} 的边段（详见 §3.4）。

\section{坐标量化与统一词表}

\subsection{标量均匀量化}

所有采样点已经过单位化（缩放到 $[-1, 1]^3$）。对每个标量分量 $x \in [-1, 1]$ 应用均匀量化：

\begin{equation}
    q(x) = \left\lfloor \frac{x + 1}{2} \cdot N_q \right\rfloor + \text{offset}, \quad N_q = 4000,\ \text{offset} = 50.
    \label{eq:quantization}
\end{equation}

量化粒度约 $5 \times 10^{-4}$，对 1\,m 尺度的零件相当于 0.5\,mm 精度，对小特征也能保留毫米级细节。

\subsection{词表组成}

数据侧的统一词表共 4\,053 个 token，按三段式布局：

\begin{itemize}
    \item \textbf{面索引段}：$[0, 49]$，共 50 个 token，用于编码每个面在当前样本中的序号；
    \item \textbf{量化坐标段}：$[50, 4049]$，共 4\,000 个 token，对应\autoref{eq:quantization} 中 $q(x)$ 的取值范围；
    \item \textbf{结构 token 段}：$\text{START}=4050$，$\text{SEP}=4051$，$\text{END}=4052$。
\end{itemize}

数据词表大小为 $50 + 4000 + 3 = 4053$。模型侧额外引入一个 PAD token（id $= 4053$）用于批内序列长度对齐，因此模型 \texttt{vocab\_size} = 4054（即数据词表大小加上 PAD，训练日志记录为 “Dataset Vocab: 4053, Model Vocab: 4054, PAD: 4053”）。

这种三段布局让自回归模型可以从 logit 分布上直接识别当前位置应该输出哪一段，避免“几何 token 与控制 token 混淆”。

\section{序列骨架：face\_block 与 edge\_block 的组织}

\subsection{完整序列结构}

一个完整样本的 token 序列以 START 起、END 止，中间由 SEP 切分为“面段”与“边段”两部分：

\begin{verbatim}
START  [face_1 block] [face_2 block] ... [face_F block]
SEP    [edge_1 block] [edge_2 block] ... [edge_E block]  END
\end{verbatim}

\begin{figure}[!htbp]
    \centering
    \includegraphics[width=0.82\textwidth]{插图/fig3-3_seq_format.pdf}
    \caption{点级 token 序列结构示意图。完整序列由 START 开始，先顺序排列若干 face\_block，再由 SEP 切换到 edge\_block 段，最后以 END 终止；每个 face\_block 含 16 个三维采样点对应的 48 个坐标 token 和 1 个面索引 token，每个 edge\_block 含 2 个邻接面索引 token 和 4 个三维采样点对应的 12 个坐标 token。}
    \label{fig:sequence-format}
\end{figure}

\subsection{face\_block：48 个坐标 token + 1 个面索引 token}

每个面段长度为 \textbf{face\_block = 49} 个 token：前 48 个 token 是当前面 16 个三维采样点按 $(x, y, z)$ 展平后的标量坐标 token，最后 1 个 token 是面索引 token。分组实现将 48 个坐标 token 按 3 个一组组织为 16 个 $[x, y, z]$ 三元组，再追加一个 $[\text{idx}, \text{PAD}, \text{PAD}]$ 的面索引三元组。

\subsection{edge\_block：2 个面索引 token + 12 个坐标 token}

每个边段长度为 \textbf{edge\_block = 14} 个 token：前 2 个 token 是该边两端邻面的面索引 token，后 12 个 token 是当前边 4 个三维采样点按 $(x, y, z)$ 展平后的标量坐标 token。\texttt{\_group\_tokens} 把 2 个面索引各组装为 $[\text{idx}, \text{PAD}, \text{PAD}]$ 三元组，再把 12 个坐标 token 分为 4 个 $[x, y, z]$ 三元组追加在后。

\subsection{三元组化分组的语义}

经过 \texttt{\_group\_tokens} 处理后，原始 1D token 序列被分组为形状 $(L, 3)$ 的张量：每一“行”是一个三元组，几何点位置上的行为 $[x, y, z]$（三个坐标 token），面索引或特殊符号位置上的行为 $[\text{token}, \text{PAD}, \text{PAD}]$。这种统一形状便于按窗口对齐输入。本文主实验的窗口大小为 1（即 \texttt{local\_window\_size = 1}），每一行作为标准 next-token 的单步预测目标；代码中也保留了 $W > 1$ 的窗口化分支以支持点级三分量联合建模（详 §4.4）。

\section{序列边界控制与特殊 token 的使用}

\subsection{START / SEP / END 的使用规则}

START 出现在序列首部，标记自回归起点；SEP 分隔面段与边段；END 标志生成终止。这三个 token 在数据词表中已经显式编码。

\subsection{PAD token 与训练侧约束}

PAD（id $= 4053$）是模型侧引入的填充 token，用于两类用途：(a) 批内不同序列长度的对齐填充，对应位置在 \texttt{attention\_mask} 中标记为 0；(b) 在面索引或特殊符号位置上的 $t_1$、$t_2$ 分量填充（仅在窗口化分支启用时使用）。训练阶段对 PAD 位置的 label 设为 $-100$，由 \texttt{nn.CrossEntropyLoss} 的 \texttt{ignore\_index} 参数跳过梯度更新；推理阶段当生成位置类别决定应该输出 PAD 时直接写入 PAD，不进入 softmax 采样。这一约束让模型只在“需要预测有效 token”的位置参与梯度。

\section{拓扑感知排序策略}

\subsection{设计目标}

在固定词表与序列骨架之外，面 / 边的排列顺序也是序列化设计的一部分。如果按任意顺序排列，序列中“面 $f$ 与其邻接边 $e$”之间的相对位置完全由随机置换决定，期望距离可达数百 token；在长序列上，自回归注意力对远距离信号的聚合相对困难，模型对局部一致性的建模也会更难。本文为面与边各设计了一种确定性排序方案，目的是\textbf{提高序列中拓扑相邻元素的位置邻近性、降低相同 B-rep 在不同次预处理下的样本差异}。

\begin{figure}[!htbp]
    \centering
    \includegraphics[width=0.82\textwidth]{插图/fig3-2_topology_sort.pdf}
    \caption{拓扑感知排序策略示意图。面排序以面—边邻接图为基础，通过深度优先遍历使拓扑相邻的面在序列中尽量靠近；边排序再结合已得到的面访问顺序，对每个面关联的边进行确定性排列，从而降低同一 B-rep 在不同预处理轮次中的序列差异。}
    \label{fig:topology-sort}
\end{figure}

\subsection{DFS 面排序}

面排序以面—边邻接图作为基础。算法从面索引最小的面出发，按深度优先方式遍历整张图：访问当前面，对其每条邻接边，找到共享这条边的另一个面，若未访问则压栈递归。最终得到一个完整的面访问序列 $\pi$。具体伪代码如\autoref{alg:dfs-face-sort} 所示。

\begin{algorithm}[htbp]
    \caption{DFS 面排序算法}
    \label{alg:dfs-face-sort}
    \begin{algorithmic}[1]
        \Require 面集合 $F$，边集合 $E$，面—边邻接表 $\text{adj}$
        \Ensure 有序面序列 $\pi$
        \State $\text{visited}[f] \gets \text{false}$，对所有 $f \in F$
        \State $\pi \gets [\,]$
        \State $\text{stack} \gets [\arg\min_f f.\text{id}]$
        \While{$\text{stack}$ 非空}
            \State $f \gets \text{pop}(\text{stack})$
            \If{$\text{visited}[f]$} \State \textbf{continue} \EndIf
            \State $\text{visited}[f] \gets \text{true}$
            \State 将 $f$ 追加到 $\pi$
            \For{每条 $e \in \text{adj}[f]$}
                \State 设 $f'$ 为与 $f$ 共享 $e$ 的另一面（若存在）
                \If{$f'$ 存在且未访问}
                    \State 将 $f'$ 压入 $\text{stack}$
                \EndIf
            \EndFor
        \EndWhile
        \State \Return $\pi$
    \end{algorithmic}
\end{algorithm}

DFS 保证拓扑邻接的面在序列中尽量靠近，并具有连通分量内的局部聚类性质。当 B-rep 有多个连通 Shell 时，DFS 在第一个 Shell 完成后自然继续从下一个未访问面开始，不需要额外切换逻辑。

\subsection{最大索引关联边排序}

在每个面内部，边按“邻接面索引最大者”优先排序。给定面 $f_k$ 的邻接边集合 $\{e_j\}$，对每条 $e_j$ 计算

\begin{equation}
    \text{key}(e_j) = \max_{f' \ne f_k,\ f' \in \text{adj}(e_j)} \pi^{-1}(f'),
    \label{eq:edge-key}
\end{equation}

其中 $\pi^{-1}(\cdot)$ 是面在 $\pi$ 中的位置；当 $e_j$ 为自由边时取 $-\infty$。然后按 key 降序排列。具体伪代码如\autoref{alg:max-idx-edge-sort} 所示。

\begin{algorithm}[htbp]
    \caption{最大索引关联边排序算法}
    \label{alg:max-idx-edge-sort}
    \begin{algorithmic}[1]
        \Require 面 $f_k$，其邻接边集合 $\text{adj}_e[f_k]$，面排序 $\pi$
        \Ensure $f_k$ 的有序边序列
        \State $\text{keys} \gets [\,]$
        \For{每条 $e \in \text{adj}_e[f_k]$}
            \State 设 $f'$ 为与 $f_k$ 共享 $e$ 的另一面（或为空）
            \State $k \gets \pi^{-1}(f')$ 若 $f'$ 存在；否则 $-\infty$
            \State 将 $(e, k)$ 追加到 $\text{keys}$
        \EndFor
        \State 按 $k$ 降序对 $\text{keys}$ 排序
        \State \Return $[e \text{ for } (e, k) \in \text{keys}]$
    \end{algorithmic}
\end{algorithm}

该策略让“跨越较远面索引”的边出现在面内部边段的前段，使自回归生成时模型可以先看到跨度较大的边，再补齐局部边，便于在序列早期把握全局拓扑骨架。

\subsection{排序策略的论证地位}

本文将上述 DFS 面排序与最大索引边排序作为序列化设计的一部分提供，目的是提高局部一致性与建模便利性。本论文未单独设计“随机排序 vs 拓扑感知排序”的受控对照实验，因此不对排序策略对最终有效率的独立贡献做定量声明；相关消融实验作为后续工作列入 §5.6。

\section{序列长度估计}

排序完成后，单个样本的 token 序列长度由面数 $F$ 与边数 $E$ 决定：每个面贡献 face\_block = 49 个 token、每条边贡献 edge\_block = 14 个 token，加上 3 个结构 token：

\begin{equation}
    L \approx 49 F + 14 E + 3.
    \label{eq:seq-length}
\end{equation}

在 DeepCAD 子集上常见样本的 $F$ 与 $E$ 数量级使得 $L$ 远小于 4\,096 上限。数据加载阶段会过滤 \texttt{max\_seq\_len = 4\,096} 以外的样本：经过过滤后，训练集得到 64\,325 个样本，验证集得到 6\,095 个样本，整体过滤量级很小（训练侧 26 个，验证侧 1 个）。
